\documentclass[12pt]{article}
\usepackage[utf8]{inputenc}
\usepackage[T1]{fontenc}
\usepackage[margin=1in]{geometry}
\usepackage{setspace}
\usepackage{xcolor}
\usepackage[colorlinks=true, urlcolor=blue!60!black, linkcolor=blue!60!black]{hyperref}
\setlength{\parindent}{2em}
\setlength{\parskip}{1.25em}
\renewcommand{\baselinestretch}{1.0}
\title{\textit{United States v. Pink}, 315 U.S. 203 (1942)}
\author{Case Brief}
\date{}
\begin{document}
\maketitle

\section{Case Metadata}
\begin{itemize}
  \item \textbf{Case Name:} \textit{United States v. Pink}
  \item \textbf{Citation:} \href{https://supreme.justia.com/cases/federal/us/315/203/}{315 U.S. 203} (\href{https://www.law.cornell.edu/supremecourt/text/315/203}{62 S. Ct. 552; 86 L. Ed. 796}) (1942)
  \item \textbf{Argued / Decided:} Argued December 15, 1941; Decided February 2, 1942
  \item \textbf{Court and Posture:} Supreme Court of the United States; on writ of certiorari to the Supreme Court of New York (after affirmance by the New York Court of Appeals, 284 N.Y. 555, 32 N.E.2d 552 (1940))
  \item \textbf{Briefing Book Placement:} Chapter 2: Separation of Powers and the Executive
  \item \textbf{Vote:} 5--2 (Justices Reed and Jackson did not participate)
  \item \textbf{Majority Author:} Justice William O. Douglas (joined by Justices Black, Frankfurter, Murphy, and Byrnes)
  \item \textbf{Concurring Author:} Justice Felix Frankfurter
  \item \textbf{Dissenting Author:} Chief Justice Harlan F. Stone (joined by Justice Owen J. Roberts)
  \item \textbf{Constitutional Provisions:} U.S. Const. art. II, \S~2, cl. 2 (Treaty Clause); art. II, \S~3 (Reception Clause / Recognition Power); art. VI, cl. 2 (Supremacy Clause); U.S. Const. amend. V (Due Process Clause)
\end{itemize}

\section{Facts and Procedural History}
The First Russian Insurance Company, organized under the laws of the former Russian Empire, opened a New York branch in 1907 and deposited substantial assets with New York state authorities to secure domestic policyholders and creditors (\href{https://www.law.cornell.edu/supremecourt/text/315/203}{315 U.S. at 210}). Following the Russian Revolution, the Soviet government issued decrees in 1918 and 1919 nationalizing all Russian insurance enterprises, purporting to extinguish their debts and confiscate their assets wherever situated.

In 1925, New York courts placed the company's New York branch into liquidation under the New York Superintendent of Insurance (eventually respondent Louis H. Pink). After paying all domestic creditors whose claims arose from New York branch business, the Superintendent held a surplus exceeding \$1 million. In 1931, the New York Court of Appeals directed the Superintendent to pay the claims of foreign creditors whose transactions had arisen abroad and to remit any remaining balance to the company's surviving directors (\href{https://supreme.justia.com/cases/federal/us/315/203/}{315 U.S. at 210--11, 226}).

Before distribution was completed, President Franklin D. Roosevelt formally recognized the Soviet Union on November 16, 1933. As an integral part of recognition, the President entered into an executive agreement---the Litvinov Assignment---without Senate ratification, under which the Soviet government assigned to the United States all amounts due it from American nationals and institutions, creating a fund to satisfy American citizens' confiscation claims against Russia (\href{https://www.law.cornell.edu/supremecourt/text/315/203}{315 U.S. at 211--13}). After an earlier federal suit was remitted to state court, the United States sued Superintendent Pink and designated foreign creditors in New York Supreme Court in 1934 to recover the entire remaining surplus. Relying on its prior decision in \textit{Moscow Fire Insurance Co. v. Bank of New York \& Trust Co.} (which the U.S. Supreme Court had affirmed by an equally divided 4--4 vote, 309 U.S. 624 (1940)), the New York courts dismissed the United States' complaint on summary judgment, holding that Soviet nationalization decrees had no extraterritorial effect over New York assets and that New York's public policy against confiscation governed. The Supreme Court granted certiorari.

\section{Issue(s) Presented}
\begin{enumerate}
  \item \textbf{Supremacy of Sole Executive Agreements:} Does a sole executive agreement concluded by the President incident to recognizing a foreign government (the Litvinov Assignment), without Senate advice and consent, carry the same supremacy over conflicting state law and public policy as an Article II treaty?
  \item \textbf{Displacement of Foreign Creditors and Fifth Amendment:} Does the Litvinov Assignment vest title in the United States to the surplus New York assets of a nationalized Russian insurer as against foreign creditors whose claims arose abroad, and does prioritizing the United States and its nationals over those foreign creditors violate the Fifth Amendment?
\end{enumerate}

\section{Holding}
\begin{enumerate}
  \item \textbf{Yes.} Extending \textit{United States v. Belmont}, 301 U.S. 324 (1937), the Court held that the President's constitutional authority to recognize foreign governments includes the implied power to settle outstanding claims via a sole executive agreement, and that such an international compact carries the same supremacy over inconsistent state law and public policy as a formal Article II treaty (\href{https://supreme.justia.com/cases/federal/us/315/203/}{315 U.S. at 229--33}).
  \item \textbf{Yes; no Fifth Amendment violation.} Independently construing Russian law, the Court held that the 1918--1919 Soviet decrees reached the company's New York surplus, that title passed to the United States under the Litvinov Assignment, and that neither New York's marshaling decree nor the Fifth Amendment bars the federal government from securing priority for itself and its nationals over foreign creditors whose claims arose abroad (\href{https://supreme.justia.com/cases/federal/us/315/203/}{315 U.S. at 226--29, 234}). Reversed and remanded.
\end{enumerate}

\section{Reasoning}
Writing for a five-Justice majority, Justice Douglas first cleared two procedural hurdles: the 4--4 affirmance in \textit{Moscow Fire} lacked precedential force, and because the federal right under the Litvinov Assignment depended on the scope of Soviet nationalization decrees, the Supreme Court could independently determine foreign law rather than deferring conclusively to New York courts (\href{https://caselaw.findlaw.com/court/us-supreme-court/315/203.html}{315 U.S. at 216--18}). Accepting an official declaration from the Soviet Commissariat for Justice, Douglas concluded that the 1918--1919 decrees were intended to confiscate the insurer's foreign assets as well as its domestic property.

Turning to the merits, Douglas extended \textit{Belmont} from a corporate depositor to competing foreign creditors. While acknowledging that aliens enjoy Fifth Amendment protections, he reasoned that the federal government is not barred by due process from preferring itself and American claimants over foreign creditors whose claims did not arise from business conducted with the New York branch (\href{https://supreme.justia.com/cases/federal/us/315/203/}{315 U.S. at 228}).

Douglas then anchored the Litvinov Assignment's preemptive force in exclusive federal authority over foreign affairs. Recognition of the Soviet Union and settlement of American confiscation claims were interdependent parts of one diplomatic bargain. Invoking \textit{United States v. Curtiss-Wright Export Corp.}, 299 U.S. 304 (1936), Douglas held that removing obstacles to recognition was a ``modest implied power'' of the President as ``sole organ'' in international relations, requiring no Senate consent (\href{https://www.law.cornell.edu/supremecourt/text/315/203}{315 U.S. at 229}). Because foreign-affairs power is vested exclusively in the national government, sole executive agreements like the Litvinov Assignment possess the same supremacy over state law as Article II treaties. Allowing New York to refuse enforcement---whether on grounds of territoriality, corporate separateness, or hostility to Soviet confiscation---would resurrect the exact diplomatic friction the President had resolved.

\section{Concurring and Dissenting Opinions}
\textbf{Concurrence (Justice Frankfurter):} Joining the Court's opinion, Justice Frankfurter added a methodological warning against importing domestic conflict-of-laws doctrines---such as ``situs,'' ``jurisdiction,'' and ``comity''---into international controversies (\href{https://caselaw.findlaw.com/court/us-supreme-court/315/203.html}{315 U.S. at 234}). Once the political branch recognized the Soviet Union and accepted the Litvinov Assignment as the instrument for resolving bilateral claims, state courts could not frustrate that national settlement through local public-policy exceptions.

\textbf{Dissent (Chief Justice Stone, joined by Justice Roberts):} Chief Justice Stone argued that the Court drastically overread both diplomatic recognition and the Litvinov Assignment (\href{https://supreme.justia.com/cases/federal/us/315/203/}{315 U.S. at 242--56}). In Stone's view, recognizing a foreign regime merely validates its acts within its own territory; it does not retroactively override New York property law or transfer to the Soviet Union (and hence to the United States) assets located in New York that New York had already marshaled for creditors. Even assuming the President could displace state law by an express executive agreement, Stone contended that the text of the Litvinov Assignment assigned only whatever rights the Soviet Union actually possessed---it contained no explicit command overriding state creditor-priority rules or extinguishing third-party claims.

\section{Key Quotations}
\begin{quote}
``Power to remove such obstacles to full recognition as settlement of claims of our nationals ... certainly is a modest implied power of the President who is the `sole organ of the federal government in the field of international relations.'\,''\\
--- Justice Douglas, grounding the Litvinov Assignment in Article II (\href{https://www.law.cornell.edu/supremecourt/text/315/203}{315 U.S. at 229})
\end{quote}

\begin{quote}
``Enforcement of such state policies would indeed tend to restore some of the precise impediments to friendly relations which the President intended to remove on inauguration of the policy of recognition of the Soviet Government.''\\
--- Justice Douglas, explaining why New York's distribution policy must yield (\href{https://caselaw.findlaw.com/court/us-supreme-court/315/203.html}{315 U.S. at 231--32})
\end{quote}

\begin{quote}
``What another country has done in the way of taking over property of its nationals, and especially of its corporations, is not a matter for judicial consideration here. Such nationals must look to their own government for any redress to which they may be entitled.''\\
--- Justice Douglas, quoting \textit{Belmont} on the act-of-state principle (\href{https://supreme.justia.com/cases/federal/us/315/203/}{315 U.S. at 226})
\end{quote}

\begin{quote}
``Legal ideas like other organisms cannot survive severance from their congenial environment. Concepts like `situs' and `jurisdiction' and `comity' summarize views evolved by the judicial process, in the absence of controlling legislation, for the settlement of domestic issues. To utilize such concepts for the solution of controversies international in nature, even though they are presented to the courts in the form of a private litigation, is to invoke a narrow and inadmissible frame of reference.''\\
--- Justice Frankfurter, concurring (\href{https://caselaw.findlaw.com/court/us-supreme-court/315/203.html}{315 U.S. at 234})
\end{quote}

\section{Doctrinal Significance and Modern Legacy (Through 2026)}
\textit{United States v. Pink} represents the high-water mark of unilateral presidential preemption of state law via sole executive agreements. Together with \textit{Belmont}, it established that a President's diplomatic claims-settlement agreement---though never submitted to the Senate under the Treaty Clause and not literally enumerated in the Supremacy Clause---overrides inconsistent state statutes, judicial decrees, and public policies (\href{https://www.law.cornell.edu/constitution-conan/article-2/section-2/clause-2/legal-effect-of-executive-agreements}{Constitution Annotated}). In the 1950s, \textit{Pink}'s breadth fueled the proposed Bricker Amendment, which sought (unsuccessfully) to subject executive agreements to congressional control.

The Supreme Court later reaffirmed \textit{Pink}'s core claims-settlement principle while progressively cabining its reach. In \textit{Dames \& Moore v. Regan}, 453 U.S. 654 (1981), the Court upheld President Carter's Algiers Accords suspending private U.S. lawsuits against Iran, relying heavily on historical congressional acquiescence alongside \textit{Belmont} and \textit{Pink}. In \href{https://supreme.justia.com/cases/federal/us/539/396/}{\textit{American Insurance Ass'n v. Garamendi}}, 539 U.S. 396 (2003), a 5--4 Court relied directly on \textit{Belmont} and \textit{Pink} to hold that President Clinton's executive agreements with Germany establishing a voluntary Holocaust-era insurance compensation mechanism preempted California's mandatory disclosure statute, even without an express preemption clause.

Through 2026, however, two later landmarks define \textit{Pink}'s outer boundaries. First, in \href{https://supreme.justia.com/cases/federal/us/552/491/}{\textit{Medell\'in v. Texas}}, 552 U.S. 491 (2008), the Court sharply distinguished \textit{Belmont}, \textit{Pink}, and \textit{Garamendi} as a ``narrow set of circumstances'' involving longstanding executive practice in settling foreign claims with congressional acquiescence, refusing to let a presidential memorandum enforce a non-self-executing treaty obligation against Texas criminal procedure. Second, in \href{https://supreme.justia.com/cases/federal/us/576/1/}{\textit{Zivotofsky v. Kerry}}, 576 U.S. 1 (2015), the Court confirmed that the President holds \emph{exclusive} power to formally recognize foreign sovereigns (citing the recognition-power lineage of \textit{Pink}), while disavowing \textit{Curtiss-Wright}'s sweeping ``sole organ'' dicta as a blanket license for unbounded executive foreign-affairs power. Thus, \textit{Pink} remains binding law for presidential recognition and foreign-claims settlements, but does not confer general executive authority to make domestic law or override federal statutes.

\section{Pedagogical Critiques and Exam Hypotheticals}
\textbf{Pedagogical Critique:} \textit{Pink} poses a classic textualist puzzle under the Supremacy Clause (Article VI, cl. 2), which lists only the ``Constitution,'' ``Laws of the United States,'' and ``Treaties''---not sole executive agreements---as supreme law of the land. Students should also press Chief Justice Stone's federalism and separation-of-powers critique: if the President can bypass the Senate's two-thirds supermajority under Article II, \S~2, cl. 2, and unilaterally divest third-party creditors' state-law property rights through a diplomatic exchange of letters, what stops the Executive from using foreign-policy agreements to rewrite domestic private law?

\textbf{``Opposite Side'' Contrast (\href{https://supreme.justia.com/cases/federal/us/552/491/}{\textit{Medell\'in v. Texas}}, 552 U.S. 491 (2008)):} Pair \textit{Pink} with \textit{Medell\'in} to map the limits of executive preemption. In \textit{Pink}, a sole executive agreement tied to formal recognition and historical claims-settlement practice displaced New York's state-law distribution of foreign insurer assets. In \textit{Medell\'in}, by contrast, President George W. Bush issued a unilateral memorandum directing state courts to give effect to an International Court of Justice judgment under a non-self-executing treaty; the Supreme Court held that the President lacked authority to preempt Texas procedural-default rules. Unlike \textit{Pink}'s narrow, historically sanctioned claims-settlement context, \textit{Medell\'in} involved an unprecedented unilateral effort to convert a non-self-executing international obligation into binding domestic law without Congress.

\textbf{Exam Hypothetical:} To normalize diplomatic relations with Nation Z after a civil war, the President signs a sole executive agreement---never submitted to the Senate or Congress---that (1) assigns all U.S.-situs assets of Nation Z's nationalized state airline to the United States to compensate American passengers, and (2) requires all U.S. states to waive their state-court statutes of limitations for tort suits brought by Nation Z's diplomats against American citizens. A state court refuses to enforce either provision. How should the Supreme Court rule? \textit{Analysis:} Provision (1) sits squarely inside \textit{Belmont}, \textit{Pink}, and \textit{Garamendi}: it is a classic claims-settlement agreement incident to formal recognition, carrying preemptive force over conflicting state asset-distribution rules. Provision (2), however, falls on the \textit{Medell\'in} side of the line: altering general state tort limitations rules for private litigation lacks the historical pedigree of claims settlement and congressional acquiescence, impermissibly turning unilateral executive diplomacy into domestic lawmaking.

\end{document}
